\documentclass[BTech]{iitmdiss}
\usepackage{titlesec}
\usepackage{times}
\usepackage{t1enc}
\usepackage{float}
\usepackage{amsmath, amssymb}
\usepackage{mathrsfs}
\usepackage{longtable}
\usepackage{lipsum} % just for dummy text- not needed for a longtable
\usepackage{amsfonts}
\usepackage{amstext}
\usepackage{threeparttable}
\usepackage{textcomp}
\usepackage{rotating}
\usepackage{lscape}
\usepackage{gensymb}
\usepackage{placeins}
\usepackage{graphicx}
\usepackage{epstopdf}
\usepackage{wrapfig}
\usepackage{tabularx}
\usepackage{subfig}
\usepackage{multirow}
\usepackage{amsmath} % easier math formulae, align, subequations \ldots
\usepackage{makecell}
\usepackage{hhline}
\usepackage{cellspace}
\usepackage{enumitem}
\usepackage{mwe}
%\usepackage{graphbox}
\usepackage{hyperref}
%\hypersetup{
%	colorlinks=true, %set true if you want colored links
%	linktoc=all,     %set to all if you want both sections and subsections linked
%	linkcolor=blue,  %choose some color if you want links to stand out
%}
%\setlength\cellspacetoplimit{5pt}
%\setlength\cellspacebottomlimit{5pt}
%\usepackage{mathtools}
\usepackage{adjustbox}
% \usepackage{ltablex} 
\usepackage{inputenc}
\usepackage[english]{babel}
\usepackage[autostyle]{csquotes}
\usepackage{color}
\usepackage{framed}
\usepackage{afterpage}
\usepackage{array}
\usepackage{multirow}
\usepackage{geometry}
\usepackage{capt-of}
\usepackage{soul}
%\usepackage{tikz}
%\usetikzlibrary{shapes.geometric, arrows}
\usepackage{booktabs}
%\usepackage{pdflscape}
\usepackage{afterpage}
%\usepackage{chktex}
%\usepackage{tabu}
%\usepackage{caption}
\usepackage{bm}
\usepackage[dvipsnames,svgnames,table]{xcolor} % use color % easier math formulae, align, subequations \ldots
%\renewcommand{\baselinestretch}{1.5}
%\graphicspath{C:\Users\nitesh\Documents\Documents\IIT-Madras\My-Research\My Thesis\Nitesh_MS_Thesis\figures}
%\graphicspath{{../}{../figures/}}
%\usepackage[latin1]{inputenc}
\usepackage{tikz}
\usetikzlibrary{shapes.geometric, arrows}
\usetikzlibrary{positioning}
%%%<
\usepackage{verbatim}

\setlength{\emergencystretch}{6em}
\begin{document}
\sloppy
\hypersetup{pageanchor=false}

%%%%%%%%%%%%%%%%%%%%%%%%%%%%%%%%%%%%%%%%%%%%%%%%%%%%%%%%%%%%%%%%%%%%%%
% Title page

\title{NeuroVox}

\author{MANEPALLI MOHAN KUMAR}

\date{APRIL 2026}
\department{ELECTRONICS AND COMMUNICATION ENGINEERING}

%\nocite{*}
\maketitle

%%%%%%%%%%%%%%%%%%%%%%%%%%%%%%%%%%%%%%%%%%%%%%%%%%%%%%%%%%%%%%%%%%%%%%
% Certificate
\certificate

\vspace*{0.5in}

\noindent This is to certify that the thesis titled {\bf NeuroVox}, submitted by {\bf MANEPALLI MOHAN KUMAR}, 
  to the Rajiv Gandhi University of Knowledge Technologies, Nuzvid, for
the award of the degree of {\bf Bachelor of Technology}, is a bona fide
record of the research work done by him under our supervision. The
contents of this thesis, in full or in parts, have not been submitted
to any other Institute or University for the award of any degree or
diploma.

\vspace*{1.5in}

\begin{singlespacing}
\hspace*{-0.25in}
\parbox{2.75in}{\raggedright
\noindent {\bf Prof.~1} \\
\noindent Mr. Vinod Babu Pusuluri \\ 
\noindent Assistant Professor \\
\noindent Dept. of ECE\\
\noindent RGUKT Nuzvid, 521 201 \\
}
\end{singlespacing}


\vspace*{0.25in}
\noindent Place: Nuzvid\\
Date:  

%%%%%%%%%%%%%%%%%%%%%%%%%%%%%%%%%%%%%%%%%%%%%%%%%%%%%%%%%%%%%%%%%%%%%%
% Acknowledgements
\acknowledgements

I would like to express my deepest and most sincere gratitude to my project guide, Mr. Vinod Babu Pusuluri, Assistant Professor in the Department of Electronics and Communication Engineering at RGUKT Nuzvid. His exceptional mentorship, unwavering support, and profound expertise in IoT, embedded systems, and edge devices have been the cornerstone of my academic journey. His continuous guidance was absolutely instrumental in bringing the edge-computing architecture of this project to life, and I am profoundly thankful for the countless hours he dedicated to shaping my practical understanding of these complex domains.

I am immensely grateful to Mr. K. Shivlal, whose comprehensive teachings in deep learning laid the foundation for the machine learning methodologies applied in this work. His instruction provided me with the vital intuition and technical skills required to navigate and implement advanced neural network architectures.

A special note of thanks goes to Prof. Phaneendra K. Yalavarthy, under whose supervision I had the privilege to work as an Engineering Intern at the Medical Imaging Group Lab (MIGLab), Indian Institute of Science (IISc). I am incredibly thankful to him and the laboratory for providing critical computational resources, notably the GPU access that was indispensable for the heavy model development and training phases of this research.

Finally, I want to thank my family and peers for their unwavering encouragement, patience, and support throughout the pursuit and completion of my B.Tech degree.

%%%%%%%%%%%%%%%%%%%%%%%%%%%%%%%%%%%%%%%%%%%%%%%%%%%%%%%%%%%%%%%%%%%%%%
% Abstract
\abstract

Neurocognitive Disorders (NCDs) are traditionally diagnosed through subjective observation only after significant cognitive decline has occurred. This thesis presents NeuroVox, a multi-modal digital biomarker platform engineered to detect pre-symptomatic signatures of cognitive decay up to two to three years earlier than clinical standards, leveraging a hardware-software co-design optimized for localized edge computing. 

The system integrates three modalities to capture subtle behavioral changes. MedGemma evaluates syntactic complexity and semantic density, HeAR detects vocal micro-tremors and speech prosody variations, and computer vision on the Clock Drawing Test (CDT) identifies early visuospatial neglect and executive dysfunction. 

To ensure strict patient data privacy, the inference engine is deployed locally on NVIDIA edge architecture, bypassing cloud exposure. These multimodal streams are fused into a unified pipeline to generate a composite cognitive score, which is clinically aligned into a Mini-Mental State Examination (MMSE) equivalent. Fortified with an automated clinical alerting system for critical thresholds and validated through high-fidelity simulations, NeuroVox demonstrates a robust, scalable, and non-invasive framework bridging the critical gap between early disease onset and timely clinical intervention.

\vspace{1cm}

\noindent \textbf{Keywords:} Neurocognitive Disorders, Digital Biomarkers, Edge Computing, Multi-modal Sensor Fusion, MedGemma, Health Acoustic Representations (HeAR), Computer Vision, Privacy-Preserving Machine Learning.

\pagebreak

%%%%%%%%%%%%%%%%%%%%%%%%%%%%%%%%%%%%%%%%%%%%%%%%%%%%%%%%%%%%%%%%%
% Table of contents etc.

\begin{singlespace}
\tableofcontents
\thispagestyle{empty}

\listoftables
\addcontentsline{toc}{chapter}{LIST OF TABLES}
\listoffigures
\addcontentsline{toc}{chapter}{LIST OF FIGURES}
\end{singlespace}


%%%%%%%%%%%%%%%%%%%%%%%%%%%%%%%%%%%%%%%%%%%%%%%%%%%%%%%%%%%%%%%%%%%%%%
% Abbreviations
\abbreviations

\noindent 
\begin{tabbing}
xxxxxxxxxxx \= xxxxxxxxxxxxxxxxxxxxxxxxxxxxxxxxxxxxxxxxxxxxxxxx \kill
\textbf{IITM}   \> Indian Institute of Technology, Madras \\
\textbf{RTFM} \> Read the Fine Manual \\
\end{tabbing}

\pagebreak

%%%%%%%%%%%%%%%%%%%%%%%%%%%%%%%%%%%%%%%%%%%%%%%%%%%%%%%%%%%%%%%%%%%%%%
% Notation

\chapter*{\centerline{NOTATION}}
\addcontentsline{toc}{chapter}{NOTATION}

\begin{singlespace}
\begin{tabbing}
xxxxxxxxxxx \= xxxxxxxxxxxxxxxxxxxxxxxxxxxxxxxxxxxxxxxxxxxxxxxx \kill
\textbf{$r$}  \> Radius, $m$ \\
\textbf{$\alpha$}  \> Angle of thesis in degrees \\
\textbf{$\beta$}   \> Flight path in degrees \\
\end{tabbing}
\end{singlespace}

\pagebreak
\clearpage

% The main text will follow from this point so set the page numbering
% to arabic from here on.
\hypersetup{pageanchor=true}
\pagenumbering{arabic}

%%%%%%%%%%%%%%%%%%%%%%%%%%%%%%%%%%%%%%%%%%%%%%%%%%%%%%%%%%%%%%%%%%%%%%
% Chapter 1: Introduction
\chapter{Introduction}
\label{chap:intro}

\section{Background}
Neurocognitive Disorders (NCDs), encompassing Alzheimer's disease and various forms of dementia, present a profound and growing global health challenge. Traditionally, the diagnosis of these degenerative conditions relies heavily on subjective clinical observations and standardized neuropsychological assessments, such as the Mini-Mental State Examination (MMSE). However, a critical limitation of the current diagnostic paradigm is its reactionary nature; clinical thresholds are typically breached only after significant, irreversible neurological decay has already occurred. By the time a patient exhibits overt symptoms, therapeutic interventions are largely palliative rather than preventative. There is an urgent medical imperative to shift the diagnostic window from late-stage observation to early, pre-symptomatic detection.

\section{Motivation}
Recent advancements in artificial intelligence and sensor technologies have introduced the concept of digital biomarkers—objective, quantifiable physiological and behavioral data that can serve as indicators of disease. The motivation for this thesis stems from the recognition that cognitive decline manifests in subtle, sub-perceptual behavioral changes months or even years before it becomes clinically apparent. Micro-fluctuations in speech prosody, slight variations in syntactic complexity, and early signs of visuospatial neglect are often undetectable to the human ear or eye during routine consultations. 

Furthermore, processing these sensitive behavioral and medical modalities raises significant data privacy concerns. Relying on cloud-based infrastructure for analyzing continuous patient data exposes vulnerable demographics to potential data breaches and latency issues. Therefore, an effective early-warning system must not only be highly sensitive to multi-modal biomarkers but must also operate within a strict, privacy-preserving computational environment.

\section{Problem Statement}
The primary challenge addressed in this research is the lack of an objective, privacy-first, and multi-modal early detection framework for NCDs. Current digital screening tools often operate in isolation—analyzing either speech, language, or motor skills—leading to high false-positive rates and a lack of holistic clinical context. Additionally, the reliance on external cloud processing for computationally intensive AI models violates modern data protection standards and introduces unacceptable inference delays. This thesis aims to bridge the critical gap between early disease onset and timely clinical intervention by engineering a comprehensive, edge-deployed diagnostic pipeline.

\section{Proposed Solution: NeuroVox}
To address these challenges, this thesis introduces \textbf{NeuroVox}, an AI-powered cognitive health monitor. NeuroVox is a hardware-software co-design platform engineered to detect the earliest pre-symptomatic signatures of cognitive decay up to two to three years before standard clinical thresholds are met. 

The system utilizes a tri-modal sensor fusion approach:
\begin{itemize}
    \item \textbf{Linguistic Analysis:} Leveraging the MedGemma model to evaluate semantic density and syntactic complexity from patient transcripts.
    \item \textbf{Acoustic Analysis:} Utilizing Google's Health Acoustic Representations (HeAR) to detect vocal micro-tremors, pause frequency, and pitch variability.
    \item \textbf{Visual-Spatial Analysis:} Employing computer vision techniques to assess patient performance on the Clock Drawing Test (CDT), identifying irregular digit placement and spatial neglect.
\end{itemize}

Crucially, to ensure absolute patient data privacy and align with regulatory compliance (such as the DPDPA 2023 and HIPAA), the entire inference engine is optimized for localized edge computing. By deploying the models directly on NVIDIA edge architecture, NeuroVox completely bypasses cloud latency and external exposure. The disparate data streams are fused to generate a composite cognitive score, which is translated into an MMSE equivalent and monitored by an automated clinical alerting system to trigger real-time physician notifications when critical thresholds are breached.

\section{Objectives of the Thesis}
The specific objectives of this research project are as follows:
\begin{enumerate}
    \item To design and implement a multi-modal digital biomarker pipeline integrating linguistic, acoustic, and visual-spatial data streams.
    \item To optimize and deploy large language and acoustic models (MedGemma and HeAR) onto resource-constrained edge computing hardware for local inference.
    \item To develop a robust sensor fusion algorithm capable of synthesizing disparate data modalities into a single, clinically actionable composite cognitive score.
    \item To establish an automated alerting protocol that translates AI-derived scores into MMSE equivalents and triggers notifications at pre-defined clinical safety thresholds.
    \item To validate the core diagnostic logic and system architecture through a high-fidelity simulation environment.
\end{enumerate}

%%%%%%%%%%%%%%%%%%%%%%%%%%%%%%%%%%%%%%%%%%%%%%%%%%%%%%%%%%%%%%%%%%%%%%
% Chapter 2: Literature Review
\chapter{Literature Review}
\label{chap:lit_review}

\section{Traditional Diagnostic Paradigms for Neurocognitive Disorders}
Historically, the diagnosis of Neurocognitive Disorders (NCDs) relies on neuropsychological evaluations such as the Mini-Mental State Examination (MMSE) and the Montreal Cognitive Assessment (MoCA). While validated, these tools are inherently subjective and possess a ceiling effect, often failing to capture the earliest stages of Mild Cognitive Impairment (MCI). By the time a clinical threshold is breached on these assessments, neuronal degradation is typically advanced.

\section{Global Epidemiological Burden of Neurocognitive Disorders}
The prevalence of Neurocognitive Disorders (NCDs), particularly Alzheimer's disease and various forms of dementia, is accelerating globally at an unprecedented rate. According to the World Health Organization (WHO) and Alzheimer's Disease International, an estimated 55 million people worldwide are currently living with dementia, a figure projected to nearly triple to 139 million by the year 2050. 

The burden of these degenerative conditions is heavily distributed across both developed nations with historically aging populations and developing nations currently experiencing rapid demographic transitions. Table \ref{tab:ncd_prevalence} outlines the estimated prevalence of NCDs across several key nations, highlighting the global necessity for scalable diagnostic tools.

\begin{table}[htbp]
\centering
\caption{Estimated Prevalence of Neurocognitive Disorders (Dementia) in Select Countries}
\label{tab:ncd_prevalence}
\renewcommand{\arraystretch}{1.3}
\begin{tabularx}{\textwidth}{l l X}
\toprule
\textbf{Country} & \textbf{Estimated Cases} & \textbf{Demographic \& Clinical Context} \\
\midrule
\textbf{China} & $\sim$15.0 million & Highest absolute number of cases globally, driven by a rapidly aging population. \\
\textbf{India} & $\sim$5.3 million & High projected growth rate of NCDs due to rising life expectancy and massive population scale. \\
\textbf{United States} & $\sim$6.7 million & Major and escalating healthcare expenditure crisis, emphasizing the need for early intervention. \\
\textbf{Japan} & $\sim$6.0 million & Highest per-capita prevalence globally and often characterized as a super-aging society. \\
\textbf{Germany} & $\sim$1.8 million & Primary driver of neurological healthcare resource utilization in the European Union. \\
\textbf{United Kingdom} & $\sim$0.9 million & High disease burden with a strong national clinical focus on digital biomarker screening. \\
\bottomrule
\end{tabularx}
\end{table}

This staggering demographic reality, particularly the rising numbers in India, underscores the immediate necessity for scalable and cost-effective early detection systems. Traditional clinical observations require significant medical infrastructure, which is unfeasible at this scale. Consequently, automated platforms like NeuroVox that use localized edge computing to bypass clinical bottlenecks represent a critical evolution in managing this growing epidemiological crisis.

\section{Digital Biomarkers in Cognitive Health}
The emergence of digital biomarkers has catalyzed a shift towards objective quantification of neurological health. 
\begin{itemize}
    \item \textbf{Linguistic Biomarkers:} Research indicates that semantic density and syntactic complexity decline years before memory loss becomes apparent. Large Language Models (LLMs) like MedGemma are increasingly utilized to map these subtle linguistic erosions.
    \item \textbf{Acoustic Biomarkers:} Tools like Google's Health Acoustic Representations (HeAR) have demonstrated efficacy in capturing sub-perceptual vocal micro-tremors and altered speech prosody, which correlate highly with early-stage dementia and Parkinson's disease.
    \item \textbf{Visual-Spatial Biomarkers:} The digitized Clock Drawing Test (CDT) analyzed via computer vision provides deep insights into visuospatial neglect and executive dysfunction, mapping motor-cognitive intent against actual execution.
\end{itemize}

\section{Edge Computing for Privacy-Preserving Healthcare AI}
A significant barrier to deploying continuous digital biomarker monitoring is data privacy. Processing sensitive audio and transcript data on centralized cloud servers introduces latency and violates stringent data protection laws (e.g., DPDPA 2023). Edge computing mitigates this by bringing the inference engine directly to the data source. However, deploying computationally heavy multi-modal models onto resource-constrained ARM-based edge devices presents a significant engineering challenge, necessitating rigorous hardware-software co-design.

%%%%%%%%%%%%%%%%%%%%%%%%%%%%%%%%%%%%%%%%%%%%%%%%%%%%%%%%%%%%%%%%%%%%%%
% Chapter 3: System Architecture and Specifications
\chapter{System Architecture and Hardware Co-Design}
\label{chap:architecture}

\section{Architectural Overview}
NeuroVox is built upon a robust hardware-software co-design philosophy. The architecture is distinctly bifurcated into two environments: a high-performance compute environment for model refinement, simulation, and sensor-fusion training; and a highly constrained, privacy-preserving edge environment for real-time patient inference. 

% Placeholder for System Architecture Diagram
\begin{figure}[htbp]
    \centering
    % Replace 'architecture_diagram.png' with your actual image file later
    \rule{0.8\textwidth}{6cm} % This creates a blank placeholder box
    \caption{High-level multi-modal architecture of the NeuroVox system, detailing the flow from data acquisition to edge inference.}
    \label{fig:sys_architecture}
\end{figure}

\section{DGX Spark Development Environment}
The development, high-fidelity simulation, and complex model fine-tuning were executed utilizing the computational resources provided by the Medical Imaging Group Lab (MIGLab) at IISc. The core backbone of this phase was the NVIDIA DGX Spark infrastructure. This environment allowed for the rapid parallel processing required to handle large medical datasets, train the multi-modal fusion algorithms, and establish the baseline for the \texttt{SIMULATION\_MODE = True} pipeline.

\begin{table}[htbp]
\centering
\caption{Hardware Specifications: NVIDIA DGX Spark (Training \& Simulation)}
\label{tab:dgx_specs}
\renewcommand{\arraystretch}{1.3}
\begin{tabularx}{\textwidth}{p{4.5cm} X}
\toprule
\textbf{Component} & \textbf{Specification} \\
\midrule
\textbf{GPU Architecture} & Multi-GPU configuration optimized for high-throughput AI workloads. \\
\textbf{Compute Capability} & Extreme parallel processing for deep learning and high-fidelity simulations. \\
\textbf{Primary Role} & Multi-modal sensor fusion training, HeAR/MedGemma weight refinement, and baseline pipeline validation. \\
\textbf{Operating System} & Ubuntu Linux (x86\_64 architecture). \\
\textbf{Workload Managed} & Processing large clinical datasets and executing complex hyperparameter optimization for the composite scoring algorithm. \\
\bottomrule
\end{tabularx}
\end{table}

\section{Edge Inference Environment: NVIDIA Jetson Orin NX}
To guarantee real-time, privacy-first patient monitoring without relying on cloud infrastructure, the final inference engine of NeuroVox is deployed on the NVIDIA Jetson Orin NX. This device was selected for its unparalleled compute-to-power-consumption ratio, making it the ideal candidate for an IoT-enabled medical device.

Deploying on this hardware required careful optimization. The models, including MedGemma and the HeAR acoustic networks, were quantized and optimized using NVIDIA TensorRT to ensure they could operate smoothly within the thermal and memory constraints of the Orin NX.

% Placeholder for Jetson Orin NX hardware photo
\begin{figure}[htbp]
    \centering
    \rule{0.6\textwidth}{5cm} % Placeholder box
    \caption{The NVIDIA Jetson Orin NX edge computing module utilized for localized, privacy-preserving inference.}
    \label{fig:jetson_hardware}
\end{figure}

\begin{table}[htbp]
\centering
\caption{Hardware Specifications: NVIDIA Jetson Orin NX (Edge Inference)}
\label{tab:jetson_specs}
\renewcommand{\arraystretch}{1.3}
\begin{tabularx}{\textwidth}{p{4.5cm} X}
\toprule
\textbf{Component} & \textbf{Specification} \\
\midrule
\textbf{AI Performance} & Up to 100 TOPS (tera operations per second). \\
\textbf{GPU} & NVIDIA Ampere architecture with 1024 CUDA cores and 32 Tensor Cores. \\
\textbf{CPU} & 8-core ARM Cortex-A78AE v8.2 64-bit CPU. \\
\textbf{Memory} & 16 GB 128-bit LPDDR5. \\
\textbf{Storage} & NVMe SSD support. \\
\textbf{Power Consumption} & Configurable between 10 W and 25 W. \\
\textbf{Primary Role} & Localized, real-time inference of linguistic, acoustic, and visual-spatial models. \\
\bottomrule
\end{tabularx}
\end{table}

\section{Software Stack and ARM Deployment Challenges}
The software ecosystem for the edge device is built upon \textbf{Ubuntu 24.04 LTS (ARM64)}. Moving from the traditional x86\_64 architecture used in the DGX environment to the ARM64 architecture of the Jetson Orin NX posed specific engineering challenges. 

\begin{itemize}
    \item \textbf{Operating System:} Ubuntu 24.04 LTS was chosen for its long-term support and native integration with the latest NVIDIA JetPack SDK.
    \item \textbf{Dependency Management:} Standard Python packages for deep learning often lack pre-compiled binaries for ARM64. Extensive environment configuration, including custom compilation of libraries using specific \texttt{CXXFLAGS} and custom CUDA 11.8 paths, was required.
    \item \textbf{Model Execution:} The core logic, encapsulating the \texttt{MedGemmaEngine}, \texttt{HeAREngine}, and \texttt{VisionEngine}, is orchestrated via Python. The models are loaded into the Jetson's unified memory space, allowing zero-copy data transfer between the CPU and the Ampere GPU, drastically reducing latency during the multi-modal fusion process.
\end{itemize}

% Placeholder for Software Stack Diagram
\begin{figure}[htbp]
    \centering
    \rule{0.8\textwidth}{6cm} % Placeholder box
    \caption{Software stack and dependency tree for the ARM64 deployment on Ubuntu 24.04 LTS.}
    \label{fig:software_stack}
\end{figure}
%%%%%%%%%%%%%%%%%%%%%%%%%%%%%%%%%%%%%%%%%%%%%%%%%%%%%%%%%%%%
% Bibliography.

\begin{singlespace}
  \nocite{*}
  \bibliography{refs}
\end{singlespace}


%%%%%%%%%%%%%%%%%%%%%%%%%%%%%%%%%%%%%%%%%%%%%%%%%%%%%%%%%%%%
% List of papers

\listofpapers

\begin{enumerate}  
\item Authors....  \newblock
 Title...
  \newblock {\em Journal}, Volume,
  Page, (year).
\end{enumerate}  

\end{document}
